% !TeX root = main.tex
\changecolours{ScoutNavy}
\shorthandoff{"}
\section[Estágios Fundamentais: Filtragem e Projeção]{Estágios Fundamentais: Filtragem, Projeção e Ordenação}

% ------------------------------------------------------------------------------
\begin{frame}[fragile]{O Estágio \texttt{\$match}: Filtragem Precoce e Índices}
  \vspace{0.1cm}
  \begin{columns}[t,onlytextwidth]
    \begin{column}{0.48\textwidth}
      \alert{Princípios de Execução}\par
      \footnotesize
      \begin{itemize}
        \item Recebe os mesmos predicados de busca utilizados no método \lstinline[language=mongosh, basicstyle=\footnotesize\ttfamily]|find()|.
        \item \textbf{Aproveitamento de Índices:} Quando colocado na primeira posição do pipeline, utiliza índices $B\text{-Tree}$ da coleção.
        \item \textbf{Regra de Performance:} Reduza o volume no Estágio 1 para poupar RAM e CPU de todos os estágios seguintes.
      \end{itemize}
    \end{column}
    \begin{column}{0.48\textwidth}
      \alert{Exemplo no E-commerce ShopGrids}\par
      \vspace{0.05cm}
      \scriptsize
      Selecionar pedidos entregues $\ge$ R\$ 200,00:
      \vspace{0.05cm}
      \begin{lstlisting}[language=mongosh, basicstyle=\scriptsize\ttfamily]
db.pedidos.aggregate([
  {
    $match: {
      status: "ENTREGUE",
      "pagamento.total_pago": { $gte: 200.0 }
    }
  }
])
      \end{lstlisting}
      \vspace{0.05cm}
      \textbf{Plano de Execução:} Com índice \lstinline[language=mongosh, basicstyle=\scriptsize\ttfamily]|{ status: 1, "pagamento.total_pago": 1 }|, o motor realiza \texttt{IXSCAN} sem ler a coleção inteira.
    \end{column}
  \end{columns}
\end{frame}

% ------------------------------------------------------------------------------
\begin{frame}[fragile]{Projeção de Documentos: \texttt{\$project} versus \texttt{\$addFields}}
  \vspace{0.1cm}
  \begin{columns}[t,onlytextwidth]
    \begin{column}{0.40\textwidth}
      \alert{Quando Usar Cada Estágio?}\par
      \footnotesize
      \begin{itemize}
        \item \textbf{\lstinline[language=mongosh, basicstyle=\footnotesize\ttfamily]|$project| (Reestruturação Completa):} Define estritamente quais campos seguem no fluxo (1 = inclui, 0 = remove). O campo \lstinline[language=mongosh, basicstyle=\footnotesize\ttfamily]|_id| é mantido por padrão e requer \lstinline[language=mongosh, basicstyle=\footnotesize\ttfamily]|_id: 0| para ser suprimido.
        \item \textbf{\lstinline[language=mongosh, basicstyle=\footnotesize\ttfamily]|$addFields| (Enriquecimento Suave):} Mantém todos os campos existentes e injeta apenas atributos calculados, evitando reescrever a projeção inteira.
        \item \textbf{Performance:} Use \lstinline[language=mongosh, basicstyle=\footnotesize\ttfamily]|$project| precoce para descartar dados pesados da memória RAM.
      \end{itemize}
    \end{column}
    \begin{column}{0.55\textwidth}
      \alert{Demonstração de Sintaxe}\par
      \vspace{0.05cm}
      \begin{lstlisting}[language=mongosh, basicstyle=\scriptsize\ttfamily]
// 1. $project: redefine a saida projetando codigo e total
{ $project: { _id: 0, codigo: 1, total_pago: "$pagamento.total_pago" } }

// 2. $addFields: adiciona desconto sem perder campos
{
  $addFields: {
    desconto_aplicado: {
      $subtract: ["$pagamento.subtotal", "$pagamento.total_pago"]
    }
  }
}
      \end{lstlisting}
    \end{column}
  \end{columns}
\end{frame}

% ------------------------------------------------------------------------------
\begin{frame}[fragile]{Ordenação e Paginação: \texttt{\$sort}, \texttt{\$limit} e \texttt{\$skip}}
  \vspace{0.1cm}
  \begin{columns}[t,onlytextwidth]
    \begin{column}{0.48\textwidth}
      \alert{Controle de Fluxo e Limite de RAM}\par
      \footnotesize
      \begin{itemize}
        \item \textbf{\lstinline[language=mongosh, basicstyle=\footnotesize\ttfamily]|$sort|:} Ordena documentos (1 para ascendente, -1 para descendente).
        \item \textbf{\lstinline[language=mongosh, basicstyle=\footnotesize\ttfamily]|$limit| e \lstinline[language=mongosh, basicstyle=\footnotesize\ttfamily]|$skip|:} Realizam a paginação determinística dos registros.
        \item \textbf{Gargalo Fatal de Memória RAM:} Sem índice, o \lstinline[language=mongosh, basicstyle=\footnotesize\ttfamily]|$sort| aloca buffer em RAM limitado estritamente a \textbf{100 MB}.
        \item \textbf{Escrita Temporária:} Para relatórios pesados em lote (batch), utilize \lstinline[language=mongosh, basicstyle=\footnotesize\ttfamily]|{ allowDiskUse: true }|.
      \end{itemize}
    \end{column}
    \begin{column}{0.48\textwidth}
      \alert{Pipeline de Paginação no MongoSH}\par
      \vspace{0.05cm}
      \begin{lstlisting}[language=mongosh, basicstyle=\scriptsize\ttfamily]
db.pedidos.aggregate([
  { $match: { status: "ENTREGUE" } },
  { $sort: { data_criacao: -1 } },
  { $skip: 10 },
  { $limit: 5 },
  {
    $project: {
      _id: 0,
      codigo: 1,
      cliente_id: 1,
      total_pago: "$pagamento.total_pago"
    }
  }
])
      \end{lstlisting}
    \end{column}
  \end{columns}
\end{frame}
